\section{Introduction}
\label{sec:introduction}

World foundation models (WFMs) have rapidly progressed from compact latent-dynamics modules~\cite{ha2018world,hafner2020dreamer} to large-scale video generation systems capable of simulating complex physical environments~\cite{nvidia2025cosmos,liu2024sora,yang2024cogvideox,gao2024vista}. By learning to predict future visual observations conditioned on actions or text, WFMs serve as \emph{world simulators} that can be leveraged for model-based planning~\cite{wu2023daydreamer}, policy evolution~\cite{du2024learning}, and data synthesis in robotic learning~\cite{zhou2024robodreamer}. The Cosmos platform~\cite{nvidia2025cosmos}, in particular, has demonstrated that diffusion-transformer (DiT) architectures~\cite{peebles2023dit} trained on internet-scale video data can produce temporally coherent, physically plausible visual rollouts, making WFMs a promising backbone for embodied intelligence.

However, robotic manipulation systems are inherently \emph{multi-view} systems. Standard configurations employ wrist-mounted and egocentric cameras simultaneously to provide complementary geometric and semantic information for policy learning~\cite{brohan2023rt1,brohan2023rt2,collaboration2023openx}. When a world model is deployed as a simulated generator for such systems, it must generate future observations across all viewpoints while maintaining strict 3D consistency: the same object must appear at geometrically compatible locations, with coherent depth and texture, across all views at every time step. Failure to maintain this consistency--manifesting as cross-view object drift, depth contradictions, or texture misalignment--directly undermines the physical plausibility of imagined trajectories and propagates errors to downstream planning and control~\cite{wu2023daydreamer,chi2023diffusionpolicy}.

Existing approaches to multi-view world modeling fall short of this requirement. Single-view WFMs such as Cosmos~\cite{nvidia2025cosmos}, CogVideoX~\cite{yang2024cogvideox}, and Vista~\cite{gao2024vista} produce high-quality temporal predictions but are architecturally restricted to one viewpoint. Methods that do handle multiple views, such as Genie~\cite{bruce2024genie} and iVideoGPT~\cite{wu2024ivideogpt}, typically concatenate tokens from different viewpoints along the sequence dimension without introducing explicit mechanisms for cross-view geometric reasoning. This ``flat concatenation'' strategy treats multi-view tokens identically to temporal tokens, leaving the model to discover cross-view correspondences implicitly from data. This implicit discovery becomes increasingly unreliable as the number of viewpoints and the complexity of the scene grow.

We trace the root cause of these failures to two fundamental deficiencies in existing multi-view approaches. First, they lack an explicit \emph{inter-view communication mechanism}. The flat concatenation strategy provides no dedicated pathway for viewpoints to exchange information; each view's tokens attend to the entire sequence without distinguishing same-view from cross-view tokens, forcing the model to discover geometric correspondences implicitly from data. As a result, each viewpoint effectively generates in isolation, with no way to coordinate predictions or resolve cross-view conflicts. Second, they lack a \emph{3D geometric prior}. Even with an inter-view communication pathway, the model has no supervision signal specifying what geometrically consistent 3D structure looks like. Without explicit geometric guidance, cross-view information exchange tends to converge on superficial shortcuts, such as matching color palettes or copying textures, rather than learning true 3D correspondences.

These two deficiencies demand two corresponding solutions: an information pathway that enables cross-view coordination, and a geometric learning signal that steers this pathway toward 3D-correct representations. Importantly, \emph{neither solution alone is sufficient}. An inter-view communication channel (Cross-View Attention) without geometric guidance allows information to flow across views but provides no guarantee that the exchanged information respects 3D geometry. In practice it learns shortcuts such as texture copying or uniform averaging. Conversely, a 3D geometric prior (Latent 3D-REPA, which aligns DiT representations with features from frozen 3D foundation models~\cite{wang2025vggt,lin2025depthanything3}) improves each view's 3D awareness independently, but without an inter-view channel, the geometric constraints cannot propagate across viewpoints and cross-view inconsistencies persist. Only when both are present does the system achieve genuine multi-view 3D consistency: the communication channel carries geometrically meaningful content, and the geometric prior is enforced coherently across all views.

Based on this analysis, we present \textbf{\methodname{}}, a framework that augments the Cosmos DiT backbone with three lightweight, modular components designed to jointly satisfy both conditions (\figref{fig:pipeline}). A shared \emph{Geometric Rotary Position Embedding} (Geo-RoPE) encodes camera ray directions and extrinsic poses into the attention mechanism via rotary position encoding~\cite{su2024rope,he2024cameractrl}, providing the common geometric reference frame. \emph{Cross-View Attention} blocks, interleaved within the DiT, provide the inter-view communication channel. \emph{Latent 3D-REPA} supplies the 3D geometric prior by aligning intermediate DiT features with 3D-aware representations from 3D foundation models via the REPA framework~\cite{yu2024repa}. Our contributions are as follows:
\begin{itemize}
	\item We identify two fundamental deficiencies in existing multi-view world models, namely the lack of inter-view communication and the absence of 3D geometric priors, and show that addressing both simultaneously is necessary and sufficient for multi-view 3D consistency.
	\item We introduce \methodname{}, which satisfies both conditions through Cross-View Attention, Latent 3D-REPA, and Geometric Rotary Position Embedding. These three plug-and-play modules are applicable to any DiT-based world model.
	\item We demonstrate that \methodname{} achieves state-of-the-art multi-view 3D consistency on robotic manipulation benchmarks, and that the improved consistency directly benefits downstream embodied tasks including model-based robotic planning and world action model fine-tuning.
\end{itemize}

The remainder of this paper is organized as follows. \secref{sec:related-work} reviews related work on world models, multi-view generation, and 3D-aware representations. \secref{sec:methods} presents the \methodname{} framework in detail. \secref{sec:conclusion} summarizes our findings and discusses future directions.
